\documentclass[11pt,a4paper]{article}
\usepackage{amsmath,amssymb,amsthm}
\usepackage[margin=1in]{geometry}
\usepackage{mathtools}

\newtheorem{theorem}{Theorem}
\newtheorem{corollary}[theorem]{Corollary}
\newtheorem{lemma}[theorem]{Lemma}
\theoremstyle{definition}
\newtheorem{remark}[theorem]{Remark}

\DeclareMathOperator{\Var}{Var}
\DeclareMathOperator{\Cov}{Cov}

\title{Appendix: Fisher Information for $\gamma$ in the\\
Contamination-Augmented 4PL IRT Model}
\author{}
\date{}

\begin{document}
\maketitle

\section{Model and Notation}\label{sec:notation}

Consider the contamination-augmented four-parameter logistic (4PL) IRT model.
For model $i \in \{1,\dots,N\}$ and benchmark item $j \in \{1,\dots,M\}$,
the probability of a correct response is
\begin{equation}\label{eq:4pl}
  p_{ij}
  = c_j + (d_j - c_j)\,\sigma(\eta_{ij}),
\end{equation}
where $\sigma(x) = (1+e^{-x})^{-1}$ is the standard logistic function and the
linear predictor is
\begin{equation}\label{eq:eta}
  \eta_{ij} = a_j\,(\theta_i - b_j) + \gamma\,\tau_{ij}.
\end{equation}
Here $\theta_i$ denotes the latent ability of model~$i$;
$(a_j, b_j, c_j, d_j)$ are the item parameters (discrimination, difficulty,
lower asymptote, upper asymptote);
$\gamma$ is the contamination effect parameter;
and $\tau_{ij}$ is the exposure signal for model~$i$ on item~$j$.
Responses are conditionally independent Bernoulli:
\begin{equation}\label{eq:bernoulli}
  y_{ij} \mid \theta, a, b, c, d, \gamma, \tau
  \;\sim\; \mathrm{Bernoulli}(p_{ij}).
\end{equation}

We derive the Fisher information for $\gamma$ treating all other parameters
$(\theta, a, b, c, d)$ and the exposure $\tau$ as fixed and known.
We denote $s_{ij} = \sigma(\eta_{ij})$, so that $p_{ij} = c_j + (d_j-c_j)\,s_{ij}$.

\section{Score Function}\label{sec:score}

The log-likelihood contribution of observation $(i,j)$ is
\begin{equation}\label{eq:loglik}
  \ell_{ij}(\gamma)
  = y_{ij}\,\log p_{ij} + (1-y_{ij})\,\log(1-p_{ij}).
\end{equation}

\paragraph{Step 1: Derivative of $p_{ij}$ with respect to $\gamma$.}

By the chain rule,
\begin{equation}\label{eq:dp_dgamma}
  \frac{\partial p_{ij}}{\partial\gamma}
  = (d_j - c_j)\,\sigma'(\eta_{ij})\,\frac{\partial\eta_{ij}}{\partial\gamma}.
\end{equation}

From~\eqref{eq:eta}, $\partial\eta_{ij}/\partial\gamma = \tau_{ij}$.
The derivative of the logistic function is
\begin{equation}\label{eq:sig_deriv}
  \sigma'(\eta) = \sigma(\eta)\bigl(1-\sigma(\eta)\bigr) = s_{ij}(1-s_{ij}).
\end{equation}

Substituting into~\eqref{eq:dp_dgamma}:
\begin{equation}\label{eq:dp_final}
  \frac{\partial p_{ij}}{\partial\gamma}
  = (d_j - c_j)\,s_{ij}(1-s_{ij})\,\tau_{ij}.
\end{equation}

\paragraph{Step 2: Score function.}

For a Bernoulli likelihood, the score with respect to a probability parameter is
\begin{equation}\label{eq:bern_score}
  \frac{\partial\ell_{ij}}{\partial p_{ij}}
  = \frac{y_{ij}}{p_{ij}} - \frac{1-y_{ij}}{1-p_{ij}}
  = \frac{y_{ij} - p_{ij}}{p_{ij}(1-p_{ij})}.
\end{equation}

Applying the chain rule:
\begin{align}
  \frac{\partial\ell_{ij}}{\partial\gamma}
  &= \frac{\partial\ell_{ij}}{\partial p_{ij}}\cdot\frac{\partial p_{ij}}{\partial\gamma}
  \notag\\
  &= \frac{y_{ij} - p_{ij}}{p_{ij}(1-p_{ij})}
     \cdot (d_j - c_j)\,s_{ij}(1-s_{ij})\,\tau_{ij}.
  \label{eq:score_ij}
\end{align}

\paragraph{Step 3: Define the effective weight.}
Let
\begin{equation}\label{eq:alpha}
  \alpha_{ij}
  \coloneqq \frac{(d_j - c_j)\,s_{ij}(1-s_{ij})}{p_{ij}(1-p_{ij})}.
\end{equation}

Then~\eqref{eq:score_ij} becomes
\begin{equation}\label{eq:score_ij_short}
  \frac{\partial\ell_{ij}}{\partial\gamma}
  = \alpha_{ij}\,\tau_{ij}\,(y_{ij} - p_{ij}),
\end{equation}
and the total score function is
\begin{equation}\label{eq:total_score}
  S(\gamma) = \frac{\partial\ell}{\partial\gamma}
  = \sum_{i=1}^{N}\sum_{j=1}^{M}
    \alpha_{ij}\,\tau_{ij}\,(y_{ij} - p_{ij}).
\end{equation}


\section{Fisher Information}\label{sec:fisher}

\begin{theorem}[Fisher Information for $\gamma$ in 4PL+$\gamma$ IRT]\label{thm:fisher}
Under the contamination-augmented 4PL model
\eqref{eq:4pl}--\eqref{eq:eta} with Bernoulli responses~\eqref{eq:bernoulli},
the Fisher information for $\gamma$
(conditional on $\theta, a, b, c, d$ and $\tau$) is
\begin{equation}\label{eq:fisher}
  I(\gamma)
  = \sum_{i=1}^{N}\sum_{j=1}^{M} w_{ij}\,\tau_{ij}^2,
\end{equation}
where
\begin{equation}\label{eq:weight}
  w_{ij}
  = \frac{\bigl[(d_j - c_j)\,s_{ij}(1-s_{ij})\bigr]^2}{p_{ij}(1-p_{ij})}.
\end{equation}
\end{theorem}

\begin{proof}
We give two equivalent derivations.

\medskip\noindent\textbf{Via the score variance.}\quad
Since $E[y_{ij}] = p_{ij}$, we have $E[S(\gamma)] = 0$.
By conditional independence of the $y_{ij}$,
\begin{align}
  I(\gamma)
  &= \Var\bigl[S(\gamma)\bigr]
  = \sum_{i,j} \alpha_{ij}^2\,\tau_{ij}^2\,\Var(y_{ij})
  \notag\\
  &= \sum_{i,j} \alpha_{ij}^2\,\tau_{ij}^2\,p_{ij}(1-p_{ij})
  \label{eq:score_var}
\end{align}
since $\Var(y_{ij}) = p_{ij}(1-p_{ij})$ for Bernoulli.
Substituting the definition of $\alpha_{ij}$ from~\eqref{eq:alpha}:
\begin{align}
  \alpha_{ij}^2\,p_{ij}(1-p_{ij})
  &= \frac{[(d_j-c_j)\,s_{ij}(1-s_{ij})]^2}{[p_{ij}(1-p_{ij})]^2}
     \cdot p_{ij}(1-p_{ij})
  \notag\\
  &= \frac{[(d_j-c_j)\,s_{ij}(1-s_{ij})]^2}{p_{ij}(1-p_{ij})}
  = w_{ij}.
  \label{eq:alpha_to_w}
\end{align}
Substituting~\eqref{eq:alpha_to_w} into~\eqref{eq:score_var} gives~\eqref{eq:fisher}.

\medskip\noindent\textbf{Via the Bernoulli information formula.}\quad
For independent Bernoulli observations with probability $p_{ij}(\gamma)$,
the Fisher information is the standard result
\begin{equation}\label{eq:bern_info}
  I(\gamma)
  = \sum_{i,j} \frac{(\partial p_{ij}/\partial\gamma)^2}{p_{ij}(1-p_{ij})}.
\end{equation}
Substituting~\eqref{eq:dp_final}:
\begin{align*}
  I(\gamma)
  &= \sum_{i,j}
     \frac{[(d_j-c_j)\,s_{ij}(1-s_{ij})\,\tau_{ij}]^2}{p_{ij}(1-p_{ij})} \\
  &= \sum_{i,j}
     \underbrace{\frac{[(d_j-c_j)\,s_{ij}(1-s_{ij})]^2}{p_{ij}(1-p_{ij})}}_{w_{ij}}
     \cdot\;\tau_{ij}^2.
     \qedhere
\end{align*}
\end{proof}


\begin{remark}[2PL simplification]\label{rem:2pl}
When $c_j = 0$ and $d_j = 1$ (the 2PL submodel), we have
$p_{ij} = s_{ij} = \sigma(\eta_{ij})$, so
\begin{equation}\label{eq:w_2pl}
  w_{ij}
  = \frac{[s_{ij}(1-s_{ij})]^2}{s_{ij}(1-s_{ij})}
  = s_{ij}(1-s_{ij})
  = p_{ij}(1-p_{ij}).
\end{equation}
Under the 2PL model the Fisher information simplifies to
\begin{equation}\label{eq:fisher_2pl}
  I(\gamma)
  = \sum_{i=1}^{N}\sum_{j=1}^{M}
    p_{ij}(1-p_{ij})\,\tau_{ij}^2.
\end{equation}
Each observation contributes most information when $p_{ij} = 1/2$
(i.e., items near the ability level) and when the exposure signal $|\tau_{ij}|$
is large.
\end{remark}


\begin{remark}[Conditional vs.\ marginal information]\label{rem:marginal}
Theorem~\ref{thm:fisher} gives the \emph{conditional} Fisher information,
treating $(\theta, a, b, c, d)$ as known.
In a Bayesian setting where all parameters are estimated jointly,
the marginal posterior variance $\Var(\gamma\mid\mathbf{y})$ will generally
exceed $1/I(\gamma)$ due to (a)~prior information and (b)~posterior correlations
between $\gamma$ and other parameters.
Therefore $1/I(\gamma)$ serves as a \emph{lower bound} on the achievable
posterior variance for~$\gamma$.
The tightness of this bound depends on the degree of confounding between
$\gamma\tau$ and the other terms in the linear predictor~$\eta_{ij}$.
\end{remark}


\section{Connection to Signal AUC}\label{sec:auc}

We now relate $I(\gamma)$ to the quality of the exposure signal~$\tau$,
as measured by its area under the ROC curve (AUC) for discriminating
contaminated from clean items.

\subsection{Bi-Normal Signal Model}\label{ssec:binormal}

Suppose the ground-truth contamination status $g_{ij}\in\{0,1\}$ satisfies
$g_{ij}\sim\mathrm{Bernoulli}(\pi)$ independently,
with $\pi$ denoting the contamination rate.
The exposure signal is
\begin{equation}\label{eq:exposure_dgp}
  \tau_{ij} = g_{ij} + \varepsilon_{ij},
  \quad \varepsilon_{ij}\sim N(0,\sigma^2),
\end{equation}
so that
$\tau_{ij}\mid g_{ij}=1 \sim N(1,\sigma^2)$ and
$\tau_{ij}\mid g_{ij}=0 \sim N(0,\sigma^2)$.

The AUC for discriminating $g=1$ from $g=0$ using $\tau$ is
\begin{equation}\label{eq:auc_formula}
  A = \Phi\!\left(\frac{\Delta}{\sigma\sqrt{2}}\right),
\end{equation}
where $\Delta = 1$ is the signal separation and
$\Phi$ is the standard normal CDF.
Solving for $\sigma$:
\begin{equation}\label{eq:sigma_from_auc}
  \sigma
  = \frac{1}{\sqrt{2}\,\Phi^{-1}(A)}.
\end{equation}


\subsection{AUC--Information Relationship}

\begin{corollary}[AUC--Information Relationship]\label{cor:auc}
Under the bi-normal signal model~\eqref{eq:exposure_dgp},
evaluating the Fisher information at $\gamma=0$ and taking the
expectation over $\tau$:
\begin{equation}\label{eq:EI_auc}
  \mathbb{E}_\tau\!\left[I(\gamma)\big|_{\gamma=0}\right]
  = \overline{W}\cdot NM\cdot
    \left[\frac{1}{2\bigl[\Phi^{-1}(A)\bigr]^2} + \pi\right],
\end{equation}
where $\overline{W} = (NM)^{-1}\sum_{i,j}w_{ij}^{(0)}$ is the average weight
evaluated at $\gamma=0$.
\end{corollary}

\begin{proof}
At $\gamma=0$, the linear predictor $\eta_{ij}=a_j(\theta_i-b_j)$ is
independent of $\tau$, so the weights $w_{ij}^{(0)}$ are deterministic
(they depend only on item and person parameters).
By Theorem~\ref{thm:fisher}:
\begin{equation}\label{eq:EI_step1}
  \mathbb{E}_\tau\!\left[I(\gamma)\big|_{\gamma=0}\right]
  = \sum_{i,j}w_{ij}^{(0)}\cdot\mathbb{E}[\tau_{ij}^2].
\end{equation}

We compute $\mathbb{E}[\tau_{ij}^2]$ under the mixture~\eqref{eq:exposure_dgp}:
\begin{align}
  \mathbb{E}[\tau^2]
  &= \pi\,\mathbb{E}[\tau^2\mid g=1] + (1-\pi)\,\mathbb{E}[\tau^2\mid g=0]
  \notag\\
  &= \pi\,(1^2+\sigma^2) + (1-\pi)\,(0^2+\sigma^2)
  \notag\\
  &= \pi + \pi\sigma^2 + (1-\pi)\sigma^2
  \notag\\
  &= \sigma^2 + \pi.
  \label{eq:Etau2}
\end{align}

Since the $\tau_{ij}$ are i.i.d.,
\begin{equation}
  \sum_{i,j}w_{ij}^{(0)}\cdot\mathbb{E}[\tau_{ij}^2]
  = \left(\sum_{i,j}w_{ij}^{(0)}\right)(\sigma^2+\pi)
  = \overline{W}\cdot NM\cdot(\sigma^2+\pi).
\end{equation}

Substituting $\sigma^2 = 1/\bigl(2[\Phi^{-1}(A)]^2\bigr)$
from~\eqref{eq:sigma_from_auc} yields~\eqref{eq:EI_auc}.
\end{proof}


\begin{remark}[Decomposition into signal and noise]\label{rem:decomp}
The second moment decomposes as
\[
  \mathbb{E}[\tau^2]
  = \underbrace{\sigma^2}_{\text{noise}}
  + \underbrace{\pi}_{\text{signal}}\!,
\]
where $\sigma^2 = 1/(2[\Phi^{-1}(A)]^2)$ decreases monotonically
from $\infty$ to $0$ as $A$ increases from $0.5$ to $1$.
Equivalently,
$\mathbb{E}[\tau^2] = \Var(\tau) + [\mathbb{E}(\tau)]^2$
where $\Var(\tau)=\sigma^2+\pi(1{-}\pi)$ and $\mathbb{E}(\tau)=\pi$.

In the regime $A\to 1$ (perfect signal, $\sigma\to 0$),
the Fisher information approaches $\overline{W}\cdot NM\cdot\pi$,
determined by the fraction of contaminated items alone.
In the regime $A\to 0.5$ (no signal, $\sigma\to\infty$),
$I(\gamma)\to\infty$ because the large noise in $\tau$ creates a
``well-conditioned design'' for estimating the coefficient $\gamma$,
even though $\gamma$ no longer measures a meaningful contamination effect.
This distinction between statistical precision and interpretability
motivates the use of AUC as a \emph{signal quality} diagnostic
rather than relying on $I(\gamma)$ alone.
\end{remark}


\begin{remark}[General bi-normal model]\label{rem:general_binormal}
For a general bi-normal model with signal separation $\Delta$
(not necessarily~$1$) and common noise variance~$\sigma^2$:
\[
  \tau\mid g{=}1 \sim N(\mu_1,\sigma^2),\quad
  \tau\mid g{=}0 \sim N(\mu_0,\sigma^2),\quad
  \Delta = \mu_1-\mu_0,
\]
the AUC is $A=\Phi(\Delta/(\sigma\sqrt{2}))$ and
\[
  \Var(\tau) = \sigma^2 + \pi(1{-}\pi)\Delta^2,
  \qquad
  \mathbb{E}[\tau^2] = \sigma^2 + \pi\mu_1^2 + (1{-}\pi)\mu_0^2.
\]
Fixing $\sigma$ and increasing $\Delta$ (hence AUC) increases the Fisher
information, which is the intuitive direction.
\end{remark}


\section{Measurement Error Attenuation}\label{sec:me}

When the exposure signal is observed with additional measurement error,
the effective Fisher information is attenuated.

\begin{corollary}[Measurement Error Attenuation]\label{cor:me}
Suppose the true exposure $\tau_{ij}$ is centered
($\mathbb{E}[\tau]=0$, $\Var(\tau)=\sigma_\tau^2$),
but we observe
\begin{equation}\label{eq:me_model}
  \tilde{e}_{ij} = \tau_{ij} + \epsilon_{ij},
  \qquad \epsilon_{ij}\sim N(0,\sigma_e^2),
\end{equation}
with $\epsilon$ independent of $\tau$.
If we substitute the posterior mean
$\hat{\tau}_{ij} = \mathbb{E}[\tau_{ij}\mid\tilde{e}_{ij}]$
for $\tau_{ij}$ in the Fisher information formula
(the analytical measurement error model), then
\begin{equation}\label{eq:me_ratio}
  \frac{\mathbb{E}\bigl[I_{\mathrm{ME}}(\gamma)\bigr]}
       {\mathbb{E}\bigl[I_{\mathrm{oracle}}(\gamma)\bigr]}
  = \frac{\sigma_\tau^2}{\sigma_\tau^2+\sigma_e^2},
\end{equation}
where expectations are over $\tau$ and $\epsilon$,
and both Fisher informations are evaluated at $\gamma=0$.
\end{corollary}

\begin{proof}
Under a Gaussian model
$\tau\sim N(0,\sigma_\tau^2)$ and
$\tilde{e}\mid\tau\sim N(\tau,\sigma_e^2)$,
the posterior distribution of $\tau$ given $\tilde{e}$ is
\begin{equation}\label{eq:posterior}
  \tau\mid\tilde{e} \;\sim\;
  N\!\left(
    \frac{\sigma_\tau^2}{\sigma_\tau^2+\sigma_e^2}\,\tilde{e},\;
    \frac{\sigma_\tau^2\,\sigma_e^2}{\sigma_\tau^2+\sigma_e^2}
  \right).
\end{equation}

The posterior mean is the linear shrinkage estimator:
\begin{equation}\label{eq:shrinkage}
  \hat{\tau}
  = \mathbb{E}[\tau\mid\tilde{e}]
  = \lambda\,\tilde{e},
  \qquad\text{where}\quad
  \lambda = \frac{\sigma_\tau^2}{\sigma_\tau^2+\sigma_e^2}.
\end{equation}

\paragraph{Step 1: Compute $\mathbb{E}[\hat\tau^2]$.}
Since $\hat\tau = \lambda\,\tilde{e}$ and $\tilde{e} = \tau+\epsilon$
with $\tau\perp\epsilon$:
\begin{align}
  \mathbb{E}[\hat\tau^2]
  &= \lambda^2\,\mathbb{E}[\tilde{e}^2]
  = \lambda^2\,(\sigma_\tau^2+\sigma_e^2)
  \notag\\
  &= \frac{\sigma_\tau^4}{(\sigma_\tau^2+\sigma_e^2)^2}
     \cdot(\sigma_\tau^2+\sigma_e^2)
  = \frac{\sigma_\tau^4}{\sigma_\tau^2+\sigma_e^2}.
  \label{eq:E_tauhat_sq}
\end{align}

\paragraph{Step 2: Compute the ratio.}
Using $I(\gamma)|_{\gamma=0} = \sum_{i,j}w_{ij}^{(0)}\,\tau_{ij}^2$
with the weights independent of $\tau$ at $\gamma=0$:
\begin{align}
  \mathbb{E}[I_{\mathrm{ME}}(\gamma)]
  &= \sum_{i,j}w_{ij}^{(0)}\,\mathbb{E}[\hat\tau_{ij}^2]
  = W_{\mathrm{total}}\cdot\frac{\sigma_\tau^4}{\sigma_\tau^2+\sigma_e^2},
  \label{eq:E_I_me}
  \\
  \mathbb{E}[I_{\mathrm{oracle}}(\gamma)]
  &= \sum_{i,j}w_{ij}^{(0)}\,\mathbb{E}[\tau_{ij}^2]
  = W_{\mathrm{total}}\cdot\sigma_\tau^2,
  \label{eq:E_I_oracle}
\end{align}
where $W_{\mathrm{total}}=\sum_{i,j}w_{ij}^{(0)}$.

Taking the ratio:
\begin{equation}
  \frac{\mathbb{E}[I_{\mathrm{ME}}]}{\mathbb{E}[I_{\mathrm{oracle}}]}
  = \frac{\sigma_\tau^4/(\sigma_\tau^2+\sigma_e^2)}{\sigma_\tau^2}
  = \frac{\sigma_\tau^2}{\sigma_\tau^2+\sigma_e^2}.
  \qedhere
\end{equation}
\end{proof}


\begin{remark}[Classical attenuation analogy]\label{rem:attenuation}
The factor $\sigma_\tau^2/(\sigma_\tau^2+\sigma_e^2)$
is the reliability ratio from classical measurement error theory.
In linear regression, using a covariate measured with additive Gaussian noise
attenuates the regression coefficient by this same factor.
Here, the attenuation applies to the Fisher information rather than the
coefficient itself, because the analytical ME model uses the posterior mean
$\hat\tau$ (which is correctly scaled via shrinkage) rather than the raw
noisy observation.
The information loss arises because $\hat\tau$ has lower variance than $\tau$:
\[
  \Var(\hat\tau) = \frac{\sigma_\tau^4}{\sigma_\tau^2+\sigma_e^2}
  < \sigma_\tau^2 = \Var(\tau).
\]
\end{remark}


\begin{remark}[Non-centered exposure]\label{rem:noncentered}
When $\tau$ is not centered ($\mathbb{E}[\tau]=\mu_\tau\ne 0$),
the posterior mean becomes
$\hat\tau = \mu_\tau + \lambda(\tilde{e}-\mu_\tau)$.
In this case,
\[
  \frac{\mathbb{E}[I_{\mathrm{ME}}]}{\mathbb{E}[I_{\mathrm{oracle}}]}
  = \frac{\sigma_\tau^4/(\sigma_\tau^2+\sigma_e^2)+\mu_\tau^2}
         {\sigma_\tau^2+\mu_\tau^2},
\]
which lies between $\sigma_\tau^2/(\sigma_\tau^2+\sigma_e^2)$ and~$1$,
approaching~$1$ as $|\mu_\tau|\to\infty$.
Intuitively, the non-zero mean of $\tau$ provides ``free'' Fisher information
that is unaffected by measurement noise.
\end{remark}


\end{document}
